\section{Related Work}
\label{sec:related_work}

Low-light 3D reconstruction remains challenging because adverse illumination
not only corrupts image quality, but also weakens the reliability of multi-view
supervision~\cite{beyond_darkness,ll_gaussian}. Existing NeRF-based and
3DGS-based methods mainly address this problem from two complementary
directions: degradation-aware appearance modeling and prior-guided
reconstruction. Some methods combine both directions, but usually still place
primary emphasis on one of them.

\textbf{Degradation-aware appearance modeling.} This line of works primarily
improves low-light 3D reconstruction by explicitly modeling degradation and
cross-view appearance inconsistency under adverse illumination. Representative
directions include camera response modeling~\cite{gaussian_dk}, per-view color
transformation and adaptive curve adjustment~\cite{cui2025luminance},
decomposition-based appearance modeling~\cite{ll_gaussian,llnerf}, explicit
noise-aware representations and progressive denoising~\cite{zhou2025lita,ll_gaussian},
and radiative imaging formulations based on emission, absorption, and
scattering~\cite{liu2025i2nerf}. In the NeRF literature,
AlethNeRF~\cite{cui2024aleth} and its follow-up studies introduce global and
local concealing fields in volume rendering to explain low-light degradation,
while GloNeRF~\cite{glonerf} employs a per-view bilateral grid after
volumetric rendering to absorb cross-view inconsistency caused by 2D
enhancement. More recently, LO-Gaussian~\cite{lo_gaussian} combines global
filtering for brightness, contrast, and saturation with local filtering for
spatially varying illumination changes and chromatic shifts, further improving
robustness under challenging low lighting conditions.

\textbf{Prior-guided reconstruction.} Beyond explicit degradation modeling,
some methods primarily improve low-light 3D reconstruction by introducing
stronger structural, geometric, or generative priors to stabilize optimization
under severely degraded observations. Representative examples include
illumination-invariant physical structure priors and structure rendering in
LITA-GS~\cite{zhou2025lita}, dense priors from learning-based MVS together
with physical constraints and diffusion priors in LL-Gaussian~\cite{ll_gaussian},
and prior-based supervision for low-light decomposition and enhancement, such
as structure-similarity and color-constancy priors in
AlethNeRF~\cite{cui2024aleth}, gray-world and smoothness priors in
LLNeRF~\cite{llnerf}, and pre-trained prior generator LLFlow for camera
parameter estimation and training supervision in
LO-Gaussian~\cite{lo_gaussian}. By injecting auxiliary guidance beyond raw
low-light observations, these methods improve robustness when direct
multi-view supervision becomes unstable.
